% Response to Reviewers: a point-by-point rebuttal letter for a journal revision.
% Edit the "Manuscript details" block below, then write one \reviewer section
% per referee with \revcomment, \response and \change. Compile with pdflatex.
\documentclass[10pt,a4paper]{article}

% ---- Page and type -------------------------------------------------------
% A ~31 pc measure keeps Plex Serif (a wide face) near 75 characters a line.
\usepackage[a4paper,top=24mm,bottom=28mm,inner=31mm,outer=31mm,
            headheight=12pt,footskip=12mm]{geometry}
\usepackage[T1]{fontenc}
\usepackage[utf8]{inputenc}
\usepackage{plex-serif}                     % IBM Plex Serif for the text
\usepackage[scale=0.9,semibold]{plex-sans}  % Plex Sans for labels; bold = semibold
\linespread{1.12}                           % Plex has a large x-height: open the leading
\usepackage[protrusion=true,expansion=true,tracking=true]{microtype}
\usepackage{needspace,xcolor,tcolorbox,booktabs,tabularx,array,enumitem,titlesec,fancyhdr,lastpage}
\tcbuselibrary{skins,breakable}
\setlength{\parindent}{0pt}
\setlength{\parskip}{0.55\baselineskip}
\frenchspacing
\widowpenalty=10000 \clubpenalty=10000

% ---- Colours: change the accent here --------------------------------------
\definecolor{accent}{HTML}{4B3A9E}     % deep blue-violet
\definecolor{ink}{HTML}{22212A}        % near-black text
\definecolor{muted}{HTML}{6E6A73}      % labels, references, footer
\definecolor{warmgrey}{HTML}{F3F0EB}   % reviewer comment tint
\definecolor{hair}{HTML}{CFCAC2}       % hairline rules
\color{ink}

\usepackage[hidelinks]{hyperref}       % comment references link without boxes
\hypersetup{pdftitle={Response to Reviewers},pdfauthor={E. Marchetti et al.}}

% ---- Manuscript details: edit these ----------------------------------------
\newcommand\journal{Journal of Mountain Hydrology}
\newcommand\manuscriptid{JMH-2026-0418}
\newcommand\revisionno{Revision 1}
\newcommand\papertitle{Snowmelt onset from Sentinel-1 backscatter across 41 Alpine
  catchments: a seven-year validation against in situ snow sensors}
\newcommand\paperauthors{Elena Marchetti, Tobias Wyss and Ana Ribeiro}
\newcommand\letterdate{17 October 2026}

% ---- Small typographic helpers -------------------------------------------
% Letterspaced sans capitals for every label, so labels never compete with text.
\newcommand\labelfont{\sffamily\bfseries\footnotesize}
\newcommand\caps[1]{{\sffamily\footnotesize\textls[90]{\MakeUppercase{#1}}}}

% Section headings: sans, accent, with a hairline under them.
\titleformat{\section}{\sffamily\bfseries\large\color{accent}}{}{0pt}{}
  [{\vspace{2pt}\color{hair}\titlerule[0.4pt]}]
\titlespacing*{\section}{0pt}{1.6\baselineskip}{0.6\baselineskip}

% Footer: manuscript id left, page of total right.
\pagestyle{fancy}\fancyhf{}
\renewcommand{\headrulewidth}{0pt}
\fancyfoot[L]{\color{muted}\sffamily\footnotesize Response to Reviewers\enspace·\enspace\manuscriptid}
\fancyfoot[R]{\color{muted}\sffamily\footnotesize \thepage\,/\,\pageref{LastPage}}

% ---- Rebuttal macros -------------------------------------------------------
% \reviewer{Note}          starts "Reviewer N" and restarts comment numbering.
% \revcomment[key]{Text}   a numbered comment box (R1.3); key lets \ref{key} link to it.
%                          (Named \revcomment because \comment is taken by the comment package.)
% \response                starts the authors' reply; plain text follows.
% \change{p.~6, ll.~2--9}{Text}  quotes revised manuscript text with its location.
\newcounter{reviewer}
\newcounter{rcomment}[reviewer]
\renewcommand\thercomment{R\arabic{reviewer}.\arabic{rcomment}}

\newcommand\reviewer[1]{%
  \needspace{14\baselineskip}% never strand a heading and its first comment at a page foot
  \refstepcounter{reviewer}%
  \section{Reviewer \arabic{reviewer}}
  {\color{muted}\small\itshape #1\par}}

\newtcolorbox{commentbox}[1]{enhanced,breakable,colback=warmgrey,colframe=warmgrey,
  boxrule=0pt,arc=1.5pt,left=10pt,right=10pt,top=7pt,bottom=8pt,
  before skip=1.1\baselineskip,after skip=0.5\baselineskip,
  fontupper=\small\color{ink},
  before upper={{\labelfont\color{accent}#1}\hspace{0.8em}}}

\newcommand\revcomment[2][]{%
  \needspace{7\baselineskip}% keep a comment and the start of its reply together
  \refstepcounter{rcomment}%
  \if\relax\detokenize{#1}\relax\else\label{#1}\fi
  \begin{commentbox}{\thercomment}#2\end{commentbox}\nopagebreak}

\newcommand\response{{\labelfont\color{accent}Response.}\enspace}

% Quoted changes: accent bar on the left, accent text, location above.
\newtcolorbox{changebox}{enhanced,breakable,frame hidden,colback=white,
  borderline west={1.6pt}{0pt}{accent},sharp corners,
  left=11pt,right=0pt,top=1pt,bottom=1pt,
  before skip=0.4\baselineskip,after skip=0.7\baselineskip,
  fontupper=\small\color{accent}}

\newcommand\change[2]{%
  \begin{changebox}%
    {\sffamily\footnotesize\color{muted}Revised manuscript, #1}\par\smallskip
    #2%
  \end{changebox}}

\begin{document}

% ---- Letterhead ----------------------------------------------------------
{\color{accent}\rule{28pt}{2.4pt}}\par\vspace{6pt}
\begin{minipage}[b]{0.58\textwidth}
  {\fontsize{22}{26}\selectfont Response to Reviewers}
\end{minipage}\hfill
\begin{minipage}[b]{0.38\textwidth}\raggedleft
  \sffamily\footnotesize\color{muted}\linespread{1.25}\selectfont
  \journal\\ Manuscript \manuscriptid\enspace·\enspace\revisionno\\ \letterdate
\end{minipage}\par
\vspace{8pt}{\color{hair}\rule{\textwidth}{0.4pt}}\par\vspace{4pt}
{\itshape \papertitle}\par\vspace{-4pt}
{\small\color{muted}\paperauthors}\par
\vspace{0.6\baselineskip}

% ---- Letter to the editor ------------------------------------------------
Dear Dr Lindqvist,

Thank you for handling our manuscript and for the opportunity to submit a
revised version. We are grateful to both reviewers for their careful reading.
Their comments led us to extend the validation from five to seven melt seasons,
to separate wet-snow detection from the onset date itself, and to report how the
method behaves under forest cover, which we had previously left implicit.

The main changes are summarised in the table below. In the point-by-point
replies that follow, each reviewer comment is shown in a tinted box, our response
follows it, and every passage quoted from the revised manuscript is marked with a
coloured bar together with its page and line numbers. A version with tracked
changes is uploaded as a supplementary file.

We hope that the revised manuscript is now suitable for publication in the
\emph{\journal}.

Yours sincerely,\\[2pt]
Elena Marchetti, on behalf of all authors\\
{\small\color{muted}Institute for Snow and Water Research, Bolzano}

% ---- Summary of changes --------------------------------------------------
\section{Summary of changes}
{\small
\renewcommand\arraystretch{1.25}
\begin{tabularx}{\textwidth}{@{}l>{\raggedright\arraybackslash}Xl@{}}
\toprule
\caps{Section} & \caps{Change} & \caps{Raised in} \\
\midrule
2.2 Data      & Validation extended to seven seasons (2019--2025) and 63 stations & \ref{c:r1-period} \\
3.1 Method    & Onset defined as three consecutive wet-snow acquisitions, not one & \ref{c:r1-onset}, \ref{c:r2-def} \\
3.3 Method    & Threshold sensitivity from $-1$ to $-4$\,dB added (new Fig.~5) & \ref{c:r2-thresh} \\
4.2 Results   & Forest-cover stratification of errors (new Table~3) & \ref{c:r1-forest} \\
5 Discussion  & Limits for wind-crusted and refrozen snowpacks stated explicitly & \ref{c:r2-limits} \\
4.4 Results   & Radar onset compared with optical snow-cover loss & \ref{c:r2-optical} \\
Figures       & Colour-blind-safe onset scale in Fig.~3; scale bar in Fig.~1 & \ref{c:r1-fig} \\
Data          & Code and onset maps archived with a DOI & \ref{c:r2-code} \\
\bottomrule
\end{tabularx}}

% ---- Reviewer 1 ------------------------------------------------------------
\reviewer{Overall positive; asks for a longer record and a forest analysis.}

\revcomment{The manuscript is clearly written and addresses a relevant problem. My
main concerns are the length of the validation period and the treatment of
forested pixels, detailed below.}

\response We thank the reviewer for this positive assessment. Both concerns are
addressed in the replies to \ref{c:r1-period} and \ref{c:r1-forest}.

\revcomment[c:r1-period]{Five melt seasons seem short to claim robustness,
especially since 2022 was an exceptionally dry year in the Alps. Could the
authors extend the record, or at least show results without 2022?}

\response We agree. Sentinel-1 data for 2019 and 2025 are now processed, which
extends the record to seven seasons and raises the number of station-years from
212 to 371. We also report all statistics with 2022 excluded; the median
absolute error changes by less than half a day.

\change{p.~5, ll.~118--123}{We use all Sentinel-1 interferometric wide-swath
acquisitions from 1~October 2018 to 30~June 2025, covering seven melt seasons and
63 automatic snow stations (371 station-years). Excluding the anomalously dry
2022 season changes the median absolute onset error from 4.1 to 3.8~days
(Table~S2).}

\revcomment[c:r1-onset]{A single wet-snow acquisition can be triggered by a
short warm spell in mid-winter. How do the authors avoid declaring melt onset
too early in such cases?}

\response This was a real weakness of the first version. Onset is now defined as
the first of three consecutive acquisitions classified as wet, which at the
six-day revisit means at least twelve days of persistent wet snow. The
mid-winter false starts the reviewer describes drop from 14\,\% to 3\,\% of
station-years.

\change{p.~7, ll.~171--175}{Melt onset is the date of the first of three
consecutive acquisitions in which backscatter falls more than 3\,dB below the
dry-snow reference. Isolated wet acquisitions are treated as transient
mid-winter melt and do not set the onset date.}

\revcomment[c:r1-forest]{Roughly a third of the study area is forested, where
C-band backscatter is dominated by the canopy. Errors should be stratified by
forest cover.}

\response We now stratify all errors by Copernicus tree-cover density in three
classes (Table~3). Below 30\,\% cover, performance is unchanged; above 60\,\%,
the median error rises to 7.9~days, and we recommend masking these pixels.

\change{p.~10, ll.~252--256}{Errors grow with canopy density
(Table~3): the median absolute error is 3.6~days below 30\,\% tree cover,
5.2~days between 30 and 60\,\%, and 7.9~days above 60\,\%. We therefore mask
pixels with more than 60\,\% cover in the published onset maps.}

\revcomment[c:r1-fig]{Figure~3 uses a rainbow colour scale for onset dates,
which is hard to read and not accessible to colour-blind readers.}

\response Thank you; Figure~3 now uses the perceptually uniform \emph{batlow}
scale, with ticks every ten days and month names on the colour bar. The same
scale is used in the supplementary onset maps.

\revcomment{Minor points: (a) the units of backscatter are missing in
Equation~2; (b) ``SWE'' is used before it is defined; (c) the station map lacks a
scale bar.}

\response All three are corrected: (a)~Equation~2 now states dB, (b)~snow water
equivalent (SWE) is defined at first use on p.~2, and (c)~a 20\,km scale bar has
been added to Figure~1.

% ---- Reviewer 2 ------------------------------------------------------------
\reviewer{Asks for a sharper definition, a sensitivity analysis and stated limits.}

\revcomment[c:r2-def]{The text sometimes uses ``wet-snow detection'' and ``melt
onset'' interchangeably. These are different quantities and should be kept
apart throughout.}

\response We fully agree and have revised the terminology across the manuscript.
Wet-snow detection now refers only to the per-acquisition classification, and
melt onset only to the seasonal date derived from it (see also \ref{c:r1-onset}).
A sentence in the introduction makes the distinction explicit.

\change{p.~2, ll.~41--44}{We distinguish two quantities: \emph{wet-snow
detection}, a classification of each acquisition, and \emph{melt onset}, a single
date per pixel and season derived from the detection series.}

\revcomment[c:r2-thresh]{The $-3$\,dB threshold is taken from earlier work at
other sites. How sensitive are the results to this choice?}

\response We added a sensitivity analysis over thresholds from $-1$ to $-4$\,dB in
0.5\,dB steps (new Fig.~5). Errors are lowest between $-2.5$ and $-3.5$\,dB and
vary by less than one day within that range, so we keep $-3$\,dB.

\change{p.~8, ll.~198--201}{Varying the threshold between $-2.5$ and
$-3.5$\,dB changes the median absolute error by at most 0.8~days (Fig.~5);
thresholds above $-2$\,dB produce frequent false onsets from surface roughness.}

\revcomment[c:r2-limits]{The discussion should state where the method is not
expected to work, for example on wind-crusted ridges or after refreezing.}

\response We added a paragraph on these limits. Refreezing resets the detection
series but not the onset date, which is the intended behaviour; wind crusts raise
the dry-snow reference and delay detection by a median of 2.4~days at the five
ridge stations.

\change{p.~13, ll.~318--322}{The method assumes a stable dry-snow reference.
On wind-exposed ridges, crust formation raises this reference and delays
detection by a median of 2.4~days; users should interpret onset dates on such
terrain as upper bounds.}

\revcomment[c:r2-code]{Will the processing code and the onset maps be made
available?}

\response Yes. The code and the seven annual onset maps are archived on Zenodo
under an open licence, and the Data Availability statement gives the DOI.

\revcomment[c:r2-optical]{The comparison with the optical snow-cover product in Section~4.4 is
interesting but rather brief. Does the radar onset precede the loss of optical
snow cover consistently?}

\response It does, by a median of 23~days across all stations, which is the
expected interval between the onset of melt and the disappearance of the
snowpack. We expanded Section~4.4 with this result and a short explanation.

\change{p.~11, ll.~281--285}{Radar melt onset preceded the date of optical
snow-cover loss at 96\,\% of station-years, by a median of 23~days (interquartile
range 15--31~days), consistent with the time a ripe snowpack needs to melt out.}

\vspace{0.6\baselineskip}
We thank both reviewers once more; their comments have made the paper more
careful and more useful.

\end{document}
